\chapter{几何朗兰兹的非正式导论}\label{ch:12}

\begin{center}
D. Gaitsgory
\end{center}

\renewcommand{\thesection}{\arabic{section}}
\begingroup\let\thefootnote\relax\footnotetext{本章依据讲习班上的报告写成，包含三节和一个附录。}\endgroup


引入这一主题所需的全部定义，我们无法在此一一给出。读者须对代数几何、数论与表示论的基本概念有一定程度的熟悉。为弥补省略的定义，文末附有若干书目建议。

\section{函数域上 $\GL_n$ 的自守形式}

\obx{1.1}{ob:1.1}
设 $\cK$ 是正特征的整体域。也就是说，$\cK$ 是有限域 $\Fq$ 上代数曲线 $X$ 的有理函数域；这里要求 $X$ 光滑、射影、几何不可约。

记 $|X|$ 为 $X$ 的（闭）点集，它等同于 $\cK$ 的赋值。对给定点 $x$，记 $\cK_x$ 为相应的局部域，$\cO_x$ 为 $\cK_x$ 的整数环，$k_x$ 为剩余域。

回忆：$\cK$ 的阿代尔环是限制积
\[
\A:=\prod_{x\in|X|}'\cK_x .
\]

\obx{1.2}{ob:1.2}
本章只讨论代数群 $\GL_n$；相应地，$\GL_n(\A)$ 是局部紧群。粗略地说，自守形式论研究商空间 $\GL_n(\cK)\backslash\allowbreak\GL_n(\A)$ 上（复值）函数构成的空间。我们只考虑光滑函数，即关于某个开紧子群不变的函数，并把这空间记为
\[
\Funct(\GL_n(\cK)\backslash\allowbreak\GL_n(\A)).
\]
按定义，阿代尔群 $\GL_n(\A)$ 通过右平移作用在 $\Funct(\GL_n(\cK)\backslash\allowbreak\GL_n(\A))$ 上。

\obx{1.3}{ob:1.3}
回忆
\[
\sO:=\prod_{x\in|X|}\cO_x\subset\A
\]
称为整阿代尔环。对应的子群 $\GL_n(\sO)\subset\GL_n(\A)$ 是极大紧子群。以下固定 $\GL_n(\A)$ 上的一个 Haar 测度 $\mu$，使 $\mu(\GL_n(\sO))=1$。

按定义，$\Funct(\GL_n(\cK)\backslash\allowbreak\GL_n(\A))$ 中 $\GL_n(\sO)$-不变的子空间，就是双商
\[
\GL_n(\cK)\backslash\allowbreak\GL_n(\A)/\GL_n(\sO)
\]
上的函数空间。然而，子空间
\[
\Funct(\GL_n(\cK)\backslash\allowbreak\GL_n(\A)/\GL_n(\sO))\subset\Funct(\GL_n(\cK)\backslash\allowbreak\GL_n(\A))
\]
在 $\GL_n(\A)$ 的作用下并不稳定。但它带有 Hecke 代数 $\sH(\GL_n(\A),\GL_n(\sO))$ 的自然作用。

\obx{1.4}{ob:1.4}
先回忆局部 Hecke 代数的定义。对点 $x\in|X|$，即 $\cK$ 的一个位，Hecke 代数 $\sH(\GL_n(\cK_x),\GL_n(\cO_x))$ 是 $\GL_n(\cK_x)$ 上紧支撑的 $\GL_n(\cO_x)$-双不变函数构成的向量空间。代数结构由卷积给出：$f_1,f_2\mapsto f_1*f_2$，其中
\[
(f_1*f_2)(g)=\int_{\GL_n(\cK_x)}f_1(g\cdot g_1)\cdot f_2(g_1^{-1})\,\mu(dg_1),
\]
而 $\mu$ 是满足 $\mu(\GL_n(\cO_x))=1$ 的 Haar 测度。

由构造，$\sH(\GL_n(\cK_x),\GL_n(\cO_x))$ 是结合代数，其单位元是 $\GL_n(\cO_x)\subset\GL_n(\cK_x)$ 的特征函数。

\begin{annote}
这里把「卷积」写成 $(f_1*f_2)(g)=\int f_1(g g_1)f_2(g_1^{-1})\,dg_1$（原书此式），与常见写法 $(f_1*f_2)(g)=\int f_1(g_1)f_2(g_1^{-1}g)\,dg_1$ 只差一个变元代换；两种次序给出的代数互为反代数，但对交换的球 Hecke 代数 $\sH(\GL_n(\cK_x),\GL_n(\cO_x))$ 二者一致，故不影响后文的交换性。单位元之所以是 $\GL_n(\cO_x)$ 的特征函数，是因为测度已归一化 $\mu(\GL_n(\cO_x))=1$；若换用未归一化测度，它只是该代数的一个幂等元而非单位元。
\par\vspace{0.15em}{\raggedleft 【来源：本书第 6 章 §1 对局部紧群卷积与光滑表示的约定；Cassels--Fr\"ohlich《Algebraic Number Theory》第 II 章（A）】\par}
\end{annote}

设 $\pi$ 是群 $\GL_n(\cK_x)$（相应地 $\GL_n(\A)$）的（复）表示，$v$ 是 $\pi$ 中的向量。回忆：若 $v$ 关于某个开紧子群不变，则称 $v$ 光滑；若 $\pi$ 的所有向量都光滑，则称 $\pi$ 光滑。以下除特别声明外，$\GL_n(\cK_x)$（相应地 $\GL_n(\A)$）的表示都假定光滑。

\begin{othm}{命题}{1.5}{prop:1.5}
设 $\pi$ 是 $G(\cK_x)$ 的表示。则 $G(\cO_x)$-不动向量构成的子空间 $\pi^{G(\cO_x)}$ 自然地是 $\sH(\GL_n(\cK_x),\GL_n(\cO_x))$-模。此外，若 $\pi$ 不可约，则 $\pi^{G(\cO_x)}$ 是不可约 $\sH(\GL_n(\cK_x),\GL_n(\cO_x))$-模，或为 $0$。
\end{othm}

\begin{proof}
对 $v\in\pi$ 及 $G(\cK_x)$ 上任一局部常值紧支撑函数 $f$，可考虑作用
\[
f\cdot v=\int_{G(\cK_x)}f(g)\cdot g(v)\,\mu(g).
\]
当 $v\in\pi^{G(\cO_x)}$ 且 $f\in\sH(\GL_n(\cK_x),\GL_n(\cO_x))$ 时，$f\cdot v\in\pi^{G(\cO_x)}$，于是得到 $\sH(\GL_n(\cK_x),\GL_n(\cO_x))$ 在 $\pi^{G(\cO_x)}$ 上的作用。

现设 $\pi$ 不可约。须证每个 $v\in\pi^{G(\cO_x)}$ 都在 $\sH(\GL_n(\cK_x),\GL_n(\cO_x))$ 的作用下生成整个空间。按假设，任一 $v'\in\pi$ 都可写成 $f'\cdot v$，其中 $f'$ 是某个局部常值紧支撑函数。然而，若 $v'\in\pi^{G(\cO_x)}$，则
\[
f'\cdot v=f\cdot v,
\]
其中 $f$ 由 $f'$ 对 $G(\cO_x)$ 从左侧平均得到。
\end{proof}

\obx{1.6}{ob:1.6}
回忆：若 $\pi^{G(\cO_x)}\neq0$，则不可约 $G(\cK_x)$-表示 $\pi$ 称为球表示。容易证明，上述对应 $\pi\mapsto\pi^{G(\cO_x)}$ 给出双射：不可约球 $G(\cK_x)$-表示的同构类，与不可约 $\sH(\GL_n(\cK_x),\GL_n(\cO_x))$-模的同构类一一对应。

Gelfand 的一个基本定理与 Satake 的一个结果（对 $\GL_n$ 两者都不难）合起来，可以非常具体地描述代数 $\sH(\GL_n(\cK_x),\GL_n(\cO_x))$。首先，已知它交换。其次，它几乎就是多项式代数。确切地说，
\[
\sH(\GL_n(\cK_x),\GL_n(\cO_x))=\mathbb{C}[T_x^1,\dots,T_x^{n-1},T_x^n,(T_x^n)^{-1}],
\]
其中 $T_x^i\in\sH(\GL_n(\cK_x),\GL_n(\cO_x))$ 是 $\GL_n(\cO_x)$-双陪集的特征函数，该双陪集的对角 $n\times n$ 矩阵代表为
\[
\diag(\underbrace{\varpi_x,\dots,\varpi_x}_{i},1,\dots,1),
\]
其中第一组内 $\varpi_x$ 重复 $i$ 次，而 $\varpi_x\in\cO_x$ 是一致化子。\footnote{事实上，须把 $T_x^i$ 乘以 $(-|k_x|^{1/2})^{i(n-i)}$ 作归一化。}

于是，每个不可约 $\sH(\GL_n(\cK_x),\GL_n(\cO_x))$-模都是一维的；从而对每个球表示 $\pi$，子空间 $\pi^{G(\cO_x)}$ 也一维。

此外，不可约 $\sH(\GL_n(\cK_x),\GL_n(\cO_x))$-模的同构类，可以等同于
\[
\Spec(\mathbb{C}[T_x^1,\dots,T_x^{n-1},T_x^n,(T_x^n)^{-1}]),
\]
后者又可等同于群 $\GL_n(\mathbb{C})$ 中半单共轭类的集合。

\begin{annote}
原书脚注 1 要求把 $T_x^i$ 乘以 $(-|k_x|^{1/2})^{i(n-i)}$ 作归一化，而 \S3.2 中 Hecke 函子又带一个平移 $[i(n-i)]$。两处并非巧合：在层-函数对应下，$\ell$-进层上的 Tate 扭 $\Ql(1/2)$ 既对应权 $1/2$，又对应导出范畴中的平移 $[1]$，因此因子 $|k_x|^{i(n-i)/2}$ 恰好对应该平移。若把 \S1 的 $T_x^i$ 与 \S3.2 的 $T_x^i$ 直接对读，须记住这处归一化，否则两边的特征值相差 $q_x^{i(n-i)/2}$。\par\vspace{0.15em}{\raggedleft 【来源：原书脚注 1 与 \S3.2（A）；层-函数对应中的 Tate 扭约定（B）】\par}
\end{annote}

\obx{1.7}{ob:1.7}
现在过渡到整体 Hecke 代数，$\sH(\GL_n(\A),\GL_n(\sO))$。按定义，
\[
\sH(\GL_n(\A),\GL_n(\sO))=\bigotimes_{x\in|X|}'\sH(\GL_n(\cK_x),\GL_n(\cO_x)),
\]
其中限制张量积的意思是：考虑形如 $\bigotimes_{x\in|X|}f_x$ 的元素的有限线性组合，使得对几乎所有的 $x$，$f_x$ 是局部 Hecke 代数 $\sH(\GL_n(\cK_x),\GL_n(\cO_x))$ 中的单位元。

注意，每个不可约 $\GL_n(\A)$-表示 $\pi$ 都是限制张量积
\[
\pi\simeq\bigotimes_{x\in|X|}'\pi_x,
\]
其中几乎所有 $\pi_x$ 是球表示。同样，若 $\pi^{\GL_n(\sO)}\neq0$，则称不可约 $\GL_n(\A)$-表示 $\pi$ 为球表示；这等价于上述张量积中所有 $\pi_x$ 都是球表示。

于是，给定 $\GL_n(\A)$ 的不可约球表示 $\pi$，可以赋给它一族半单共轭类 $\{\gamma_x\}$，$x\in|X|$，$\gamma_x\in\GL_n(\mathbb{C})$。

\obx{1.8}{ob:1.8}
现在回到自守的情形，即商空间 $\GL_n(\cK)\backslash\allowbreak\GL_n(\A)$。回忆：上述商空间上的函数 $f$ 称为尖点的，若对所有 $g\in\GL_n(\A)$ 及所有子群 $U$ 都有
\[
\int_{U(\cK)\backslash\allowbreak U(\A)}f(u\cdot g)\,\mu(u)=0,
\]
其中 $U$ 取 $\GL_n$ 的标准抛物子群的幂幺根。

记尖点函数构成的子空间为 $\Funct_{\cusp}(\GL_n(\cK)\backslash\allowbreak\GL_n(\A))$。按定义，它关于 $\GL_n(\A)$ 的作用不变。可以证明，$\Funct_{\cusp}(\GL_n(\cK)\backslash\allowbreak\GL_n(\A))$ 中每个元素都自动紧支撑。

回忆：若
\[
\Hom_{\GL_n(\A)}(\pi,\Funct_{\cusp}(\GL_n(\cK)\backslash\allowbreak\GL_n(\A)))\neq0,
\]
则不可约表示 $\pi$ 称为尖点自守表示。

一个重要结果（多重度 $1$ 定理）说：上述 Hom 至多一维。此外，整个空间 $\Funct_{\cusp}(\GL_n(\cK)\backslash\allowbreak\GL_n(\A))$ 分解为不可约（从而是尖点自守）表示的直和，而不是直积分。

下文研究 $\GL_n(\A)$ 的球尖点自守表示。由上述讨论，我们的问题等价于描述空间 $\Funct_{\cusp}(\GL_n(\cK)\backslash\allowbreak\GL_n(\A)/\GL_n(\sO))$ 的分解——分解为算子 $T_x^i$（$i=1,\dots,n$，$x\in|X|$）的公共特征子空间之直和。我们已经知道这样的分解确实存在，且谱是单的：即每一组对应于
\[
\{\gamma_x\in\GL_n(\mathbb{C})^{ss}/\mathrm{Conj}\},\quad\forall x\in|X|
\]
的特征值，重数至多为 $1$。

另一个重要观察是：$\GL_n(\A)$ 在 $\Funct_{\cusp}(\GL_n(\cK)\backslash\allowbreak\GL_n(\A))$ 上的作用的调和分析，完全是代数的。这意味着，所有函数取值所在的域 $\mathbb{C}$ 的拓扑无关紧要。特别地，可以把 $\mathbb{C}$ 换成任意特征 $0$ 的域；为使特征空间分解存在，取它为代数闭域。我们的选择是 $\Ql$，其中 $(\ell,q)=1$。

\obx{1.9}{ob:1.9}
设 $\sigma$ 是 $\cK$ 的 Galois 群的 $n$ 维不可约 $\Ql$-表示。假定 $\sigma$ 处处无分歧，即它可以看作 \'etale 基本群 $\pi_1(X)\to\GL_n(\Ql)$ 的表示。按定义，这样的表示等同于 $X$ 上的 lisse $\ell$-进层（又称局部系统），记为 $E_\sigma$。

对每个点 $x\in|X|$，有确定的共轭类 $\mathrm{Fr}_x$（对应 $\pi_1(X)$ 中的 Frobenius 元）。记 $\sigma_x\in\GL_n(\Ql)^{ss}/\mathrm{Conj}$ 为 $\sigma(\mathrm{Fr}_x)$ 的半单部分。（已知：若 $\sigma$ 不可约，则族 $(\sigma_x,\,x\in|X|)$ 唯一决定 $\sigma$。）就局部系统 $E_\sigma$ 而言，$\sigma_x$ 对应 $\mathrm{Gal}(\overline{k}_x/k_x)$ 中 Frobenius 元在 $E_\sigma$ 于 $x$ 处纤维上的作用。

现在给出关键定义：设 $\sigma$ 是 $n$ 维无分歧不可约 Galois 表示，$\pi$ 是 $\GL_n(\A)$ 的尖点自守表示。若对一切 $x\in|X|$ 有
\[
\sigma_x=\gamma_x,
\]
则称 $\sigma$ 在朗兰兹的意义下对应于 $\pi$；其中 $\gamma_x\in\GL_n(\Ql)^{ss}/\mathrm{Conj}$ 是
\[
\Spec(\Ql[T_x^1,\dots,T_x^{n-1},T_x^n,(T_x^n)^{-1}])
\]
中附着于 $\pi_x$ 的相应点。

朗兰兹猜想（现已是 Lafforgue 的定理）断言：由上述对应关系决定，$\sigma$ 的集合与 $\pi$ 的集合之间存在双射。

\begin{annote}
这里把「处处无分歧」等同于「来自 $\pi_1(X)$ 的表示」，用的是曲线 $X$ 上的标准等价：$X$ 的有限 \'etale 覆盖与 $\pi_1(X)$ 的开子群一一对应，故无分歧表示就是 $\pi_1(X)$ 的表示，也就是 $X$ 上的 lisse 层；而对函数域，$\pi_1(X)$ 又可等同于 $\cK$ 的 Galois 群的最大无分歧商。之所以只取 $\sigma(\mathrm{Fr}_x)$ 的半单部分，是因为 $\mathrm{Fr}_x$ 本身只定义到共轭，且 $\sigma(\mathrm{Fr}_x)$ 的幂幺部分依赖 $x$ 的选取；只有半单部分是共轭类的不变量，才能作 Satake 参数去匹配 $\gamma_x$。\par\vspace{0.15em}{\raggedleft 【来源：Milne, \emph{\'Etale Cohomology}（A）；半单部分的动机见 \S1.6 与 \S3.4（I）】\par}
\end{annote}

\section{用向量丛解释}

\obx{2.1}{ob:2.1}
对我们来说，一个关键的观察是：集合 $\GL_n(\cK)\backslash\allowbreak\GL_n(\A)/\GL_n(\sO)$ 可以改用几何语言重述。

记 $\Bun_n$ 为曲线 $X$ 上秩 $n$ 向量丛的同构类集合。我们断言：$\Bun_n$ 与 $\GL_n(\cK)\backslash\allowbreak\GL_n(\A)/\GL_n(\sO)$ 之间存在自然双射。

事实上，给定向量丛 $M$，它在 Zariski 拓扑下局部同构于平凡丛 $\cO_X^{\oplus n}$。取它在 $X$ 的一般点处的平凡化 $M\otimes_{\cO_X}\cK\simeq\cK^n$，并在各个局部位处取平凡化 $M\otimes_{\cO_X}\cO_x\simeq\cO_x^n$。

对每个 $x\in|X|$ 复合这些同构，得到
\[
\cK_x^n\simeq\cO_x^n\otimes_{\cO_x}\cK_x\simeq M\otimes_{\cO_X}\cK_x
\simeq M\otimes_{\cO_X}\cK\otimes_{\cK}\cK_x\simeq\cK^n\otimes_{\cK}\cK_x\simeq\cK_x^n,
\]
即 $\GL_n(\A)$ 的一个元素。改变上述平凡化，效果是在左边乘 $\GL_n(\cK)$ 的元素，在右边乘 $\prod_{x\in|X|}\GL_n(\cO_x)=\GL_n(\sO)$ 的元素。

\begin{annote}
这里看似要作无穷多个选择：一处一般点的平凡化，加上每个位处的局部平凡化。实际上改变某个 $x$ 处的局部平凡化，相当于在 $\GL_n(\cK_x)$ 中右乘；所有位合起来，正是右乘 $\GL_n(\sO)$。改变一般点的平凡化则相当于左乘 $\GL_n(\cK)$。于是双射只依赖向量丛的同构类。另须注意，此处 $\Bun_n$ 还只是一个集合（同构类的集合）；\S3.1 才把它升级为模叠 $\sBun_n$，届时「向量丛的同构类」被「向量丛的范畴」取代。\par\vspace{0.15em}{\raggedleft 【来源：原书 \S2.1、\S3.1（A）】\par}
\end{annote}

\obx{2.2}{ob:2.2}
\emph{例}\quad 当 $n=1$ 时，$\GL_1(\A)/\GL_1(\sO)$ 是 $X$ 上的除子集合。因此双商 $\GL_1(\cK)\backslash\allowbreak\GL_1(\A)/\GL_1(\sO)$ 是除子类（即除子模主除子）的集合。上述双射重新建立了熟知的事实：除子类群等同于线丛的同构类群。

\obx{2.3}{ob:2.3}
现在说明如何在几何语言中看待 Hecke 算子。固定点 $x\in X$ 与整数 $i$，$1\le i\le n$。

记 $H_x^i$ 为下述数据的同构类集合：$(M,M',\beta)$，其中 $M$ 与 $M'$ 是 $X$ 上秩 $n$ 向量丛，$\beta$ 是 $M$ 到 $M'$ 的单射（视二者为 $X$ 上的凝聚层），使得商 $M'/M$ 长度为 $i$，且概形论地支撑在 $x$ 处。（后一条件等价于说 $M'/M$ 非典范地同构于 $k_x^i$，其中 $k_x$ 是 $x$ 处的剩余域。）

有两个自然映射，记为 $h^{\leftarrow}$ 与 $h^{\rightarrow}$，从 $H_x^i$ 到 $\Bun_n$，分别把上述三元组送到 $M$ 与 $M'$。

\begin{othm}{引理}{2.4}{lem:2.4}
给定 $M\in\Bun_n$，$h^{\leftarrow}$ 在其上的纤维可典范地等同于 $k_x$ 上向量空间 $M_x$ 中 $i$ 维子空间的集合。换言之，$(h^{\leftarrow})^{-1}(M)$ 是 Grassmann 簇 $\Gr^i(M_x)$ 的 $k_x$-点集。
\end{othm}

对 $h^{\rightarrow}$ 同样成立；此时得到 $(h^{\rightarrow})^{-1}(M')$ 是 $\Gr^{n-i}(M'_x)$ 的 $k_x$-点空间。

\begin{proof}
设 $(M,M',\beta)$ 如上。$\beta$ 虽是凝聚层的单射，取 $x$ 处的纤维后却不再单射，而得到正合列
\[
0\to\ker(\beta|_x)\to M_x\to M'_x\to(M'/M)_x\to0.
\]
特别地，$\ker(\beta|_x)$ 是 $M_x$ 的 $i$ 维子空间。

反之，给定子空间 $V\subset M_x$，定义 $M'$ 为 $M(x)$ 的子层，其局部截面形如 $t^{-1}\cdot m$，其中 $t$ 是 $X$ 在 $x$ 处的局部参数，$m$ 是 $M$ 的局部截面且在 $x$ 处的值属于 $V$。
\end{proof}

\begin{annote}
$M(x)$ 是原书的简写，指把 $M$ 沿 $x$ 扭转一次，即 $M\otimes_{\cO_X}\cO_X(x)$。它比 $M$ 多出允许在 $x$ 处有一阶极点的截面，正是形如 $t^{-1}m$ 的那些。要求 $m$ 在 $x$ 处的值落在 $V$ 内，就限定了 $M'$ 在 $x$ 处纤维的大小；这一构造的逆给出 $\ker(\beta|_x)=V$，从而正合列中的 $\ker$ 恰好跑遍 $M_x$ 的 $i$ 维子空间。\par\vspace{0.15em}{\raggedleft 【来源：原书 \S2.4（A）；$M(x)$ 的常规含义（B）】\par}
\end{annote}

\obx{2.5}{ob:2.5}
注意，由于 $k_x$ 是有限域，$h^{\leftarrow}$ 与 $h^{\rightarrow}$ 的纤维都是有限集。

考虑从空间 $\Funct(\Bun_n)$ 到自身的算子
\[
f\mapsto(h^{\rightarrow})_!((h^{\leftarrow})^*(f)).
\]
（上标 $*$ 表示沿指定态射的函数拉回，下标 $!$ 表示沿纤维求和；可与附录中的记号对照。）
\[
\Bun_n\xleftarrow{h^{\leftarrow}}H_x^i\xrightarrow{h^{\rightarrow}}\Bun_n.
\]

\begin{othm}{命题}{2.6}{prop:2.6}
上述算子等于 $T_x^i(f)$。
\end{othm}

\obx{2.7}{ob:2.7}
\emph{插叙}\quad 我们指出与模形式情形的如下类比。设 $N$ 是自然数，$p$ 是与 $N$ 互素的素数。我们将研究模曲线 $X(N)$ 上的函数（更确切地说，某些丛的截面）（它与本文的曲线 $X$ 无关）。

按定义，$X(N)$ 参数化对 $(E,\alpha)$，其中 $E$ 是（$\mathbb{C}$ 上的）椭圆曲线，$\alpha$ 是 $E^N$ 与 $(\mathbb{Z}/N\mathbb{Z})^2$ 之间的同构（这里 $E^N$ 是 $E$ 上 $N$ 阶点构成的群）。

考虑 $p$ 处的 Hecke 对应。它是一条曲线，记为 $H_p$，参数化五元组 $(E,\alpha,E',\alpha',\beta)$，其中 $(E,\alpha)$ 与 $(E',\alpha')$ 是 $X(N)$ 的两个点，$\beta$ 是 $E\to E'$ 的同源，且 $\ker(\beta)$ 的阶为 $p$。由于 $N$ 与 $p$ 互素，$\beta$ 诱导同构 $E^N\to E'^N$；我们要求它与两个群到 $(\mathbb{Z}/N\mathbb{Z})^2$ 的等同相容。

和前面一样，有自然映射
\[
X(N)\xleftarrow{h^{\leftarrow}}H_p\xrightarrow{h^{\rightarrow}}X(N),
\]
它们定义算子 $f\mapsto(h^{\rightarrow})_!((h^{\leftarrow})^*(f))$，与作用在模形式上的 Hecke 算子一致。

\obx{2.8}{ob:2.8}
回到几何情形，用几何语言写出尖点条件。设 $n=n_1+n_2$ 是分拆。考虑短正合列
\[
0\to M_1\to M\to M_2\to 0
\]
的同构类集合 $\Fl^n_{n_1,n_2}$，其中 $M_1$ 与 $M_2$ 分别是秩 $n_1$ 与 $n_2$ 的向量丛。

有自然映射 $p$ 与 $q$
\[
\Bun_n\xleftarrow{p}\Fl^n_{n_1,n_2}\xrightarrow{q}\Bun_{n_1}\times\Bun_{n_2},
\]
分别把上述短正合列送到 $M$ 与 $(M_1,M_2)$。

定义常项算子 $r^n_{n_1,n_2}:\Funct(\Bun_n)\to\Funct(\Bun_{n_1}\times\Bun_{n_2})$ 为
\[
f\mapsto q_!(p^*(f)).
\]

现在容易看出：函数 $f$ 是尖点的，当且仅当对所有 $n_1+n_2=n$、$n_1,n_2>0$ 都有 \mbox{$r^n_{n_1,n_2}(f)=0$}。

\section{从函数到层}

\obx{3.1}{ob:3.1}
下一步，把 $\Bun_n$ 不只看成集合，而看成定义在 $\Fq$ 上代数簇的有理点集，并考虑其上的 $\ell$-进层（$\ell$-进层的简述见 \S4）。

一个小麻烦是：$\Bun_n$ 背后的几何对象不是簇，而是代数叠（参看 [L-MB]），记为 $\sBun_n$。叠与概形的主要差别如下。

按 Yoneda，范畴的对象（此处即概形 $\mathcal{Y}$）由函子 $\mathrm{Schemes}\Rightarrow\mathrm{Sets}$，$S\mapsto\Hom(S,\mathcal{Y})$ 完全决定。

在 $\sBun_n$ 的情形，上述 Hom 确实有定义，但不再是集合，而是范畴。这一点在讨论 $\Bun_n$ 时已经可见：考虑 $X$ 上向量丛的范畴，比考虑向量丛同构类的集合自然得多。

一般地，对概形 $S$，定义 $\Hom(S,\sBun_n)$ 为范畴，其对象是 $X\times S$ 上秩 $n$ 的丛，态射是向量丛之间的同构。这样定义的 $\sBun_n$ 确实是代数叠。

然而，从 $\ell$-进层的观点看，叠的行为如同代数簇。特别地，若 $\mathcal{Y}$ 是叠，范畴 $D^b(\mathcal{Y})$ 有意义（参看 \S4）。此外，若 $Y$ 是范畴 $\Hom(\Spec(\Fq),\mathcal{Y})$ 中对象的同构类集合，则 $D^b(\mathcal{Y})$ 的每个对象 $\Fc$ 都定义 $Y$ 上的函数 $f$，且层-函数对应的全部性质都成立。

特别地，范畴 $D^b(\sBun_n)$ 有意义，且每个对象 $\Fc\in D^b(\sBun_n)$ 都给出 $\Bun_n$ 上的函数 $f$。

\begin{annote}
若只记住「叠是取值到群胚的函子」，上面的说法会显得突兀：既然 $\Hom(S,\sBun_n)$ 是范畴而非集合，为什么还能谈 $D^b(\sBun_n)$？关键在于 $\ell$-进层只看见叠的「局部」信息：叠有光滑覆盖，而层的拉回沿光滑覆盖是下降的，于是可以像在簇上一样定义可构造层及其导出范畴。代价是：$D^b$ 中的对象不再有「逐点取值」的意义，只有沿 $\Fq$-点上的纤维取迹才得到函数。\par\vspace{0.15em}{\raggedleft 【来源：Laumon--Moret-Bailly, \emph{Champs alg\'ebriques}（A）；叠上 \'etale 层的一般约定（B）】\par}
\end{annote}

\obx{3.2}{ob:3.2}
现在讨论 Hecke 算子在这个几何情境中变成什么。

先固定 $x\in|X|$。考虑 Hecke 叠 $\sH_x^i$：$\Hom(S,\sH_x^i)$ 是三元组 $(M,M',\beta)$ 的范畴，其中 $M$ 与 $M'$ 是 $X\times S$ 上秩 $n$ 向量丛，$\beta$ 是凝聚层的嵌入 $M\hookrightarrow M'$，使得 $M'/M$ 概形论地支撑在 $x\times S\subset X\times S$ 上，且是该子概形上秩 $i$ 的向量丛。

$\sH_x^i$ 的 $\Fq$-点的同构类集合，就是我们原来的 $H_x^i$。和前面一样，有映射 $h^{\leftarrow},h^{\rightarrow}:\sH_x^i\to\sBun_n$。但现在对 $\sBun_n$ 的给定点，即向量丛 $M$，原像 $(h^{\leftarrow})^{-1}\subset\sH_x^i$ 不再是集合，而是同构于 $\Gr^i(M_x)$ 的概形。

引入 Hecke 函子 $T_x^i:D^b(\sBun_n)\to D^b(\sBun_n)$：
\[
\Fc\mapsto(h^{\rightarrow})_!((h^{\leftarrow})^*(\Fc))[i(n-i)],
\]
其中括号内的表达式表示上同调平移。

\obx{3.3}{ob:3.3}
不过，让点 $x$ 沿曲线移动会更好。引入「大」Hecke 叠 $\sH^i$ 如下。

$\Hom(S,\sH^i)$ 是四元组 $(\phi,M,M',\beta)$ 的范畴，其中 $M,M'$ 同前，$\phi$ 是映射 $S\to X$，$\beta$ 仍是凝聚层的嵌入 $M\hookrightarrow M'$，使得商支撑在图像 $\Gamma_\phi\subset X\times S$ 上，且是其上秩 $i$ 的向量丛。

除原有的投影 $h^{\leftarrow},h^{\rightarrow}:\sH^i\to\sBun_n$ 外，现在还有自然映射 $s:\sH^i\to X$，把上述四元组 $(\phi,M,M',\beta)$ 送到 $\phi$。对固定 $x\in X$，有 $s^{-1}(x)\simeq\sH_x^i$。

定义 Hecke 函子 $T^i:D^b(\sBun_n)\to D^b(X\times\sBun_n)$ 为
\[
\Fc\mapsto(s\times h^{\rightarrow})_!((h^{\leftarrow})^*(\Fc))[i(n-i)].
\]

\obx{3.4}{ob:3.4}
现在可以给出 Hecke 本征层的定义。我们的定义稍嫌粗糙（对任意群都适用的更好定义见 [BD]）。

设 $\sigma:\pi_1(X)\to\GL_n(\Ql)$ 是表示；回忆 $E_\sigma$ 是相应的 $X$ 上 $n$ 维局部系统。

若对每个 $i=1,\dots,n$ 都存在同构
\[
T^i(\Fc_\sigma)\simeq\Lambda^i(E_\sigma)\boxtimes\Fc_\sigma,
\]
则称对象 $\Fc_\sigma\in D^b(\sBun_n)$ 为关于 $\sigma$ 的 Hecke 本征层；其中 $\Lambda^i$ 表示 $i$ 次外幂。

若基域 $k=\Fq$，且 $\Fc_\sigma$ 是关于 $\sigma$ 的 Hecke 本征层，则其对应的 $\GL_n(\cK)\backslash\allowbreak\GL_n(\A)/\GL_n(\sO)$ 上的函数，是 Hecke 算子的公共特征函数，特征值由 $\sigma$ 给出，如 \S1.9 所述。

几何朗兰兹猜想（由 [L] 在 Drinfeld 原始论文之后提出；[Dr] 完整处理了 $\GL_2$ 的情形）预测：对每个如上所述的不可约 $\sigma$，确实存在 Hecke 本征层 $\Fc_\sigma$，而且它是尖点的（即 \S2.8 所定义的函子 $r^n_{n_1,n_2}$ 在 $\Fc_\sigma$ 上为零）。

\begin{annote}
原书自称上述定义「稍嫌粗糙」，并在 [BD] 中指出对任意群都适用的更好定义。差别大致在于：这里只对 $\GL_n$ 陈述，且把本征性直接写成外幂 $\Lambda^i(E_\sigma)$ 的显式同构。对一般约化群 $G$，Hecke 函子作用在 $\Bun_G$ 上，本征层的「特征值」应取对偶群 $\check{G}$ 的局部系统，而外幂 $\Lambda^i$ 只是 $G=\GL_n$ 时标准表示的外幂。\par\vspace{0.15em}{\raggedleft 【来源：原书 \S3.4 与 [BD]（A）】\par}
\end{annote}

\obx{3.5}{ob:3.5}
最后讨论一个问题：为什么把讨论限制在球（即无分歧）情形？

理由如下：如前所述，球尖点自守表示对应 $\GL_n(\cK)\backslash\allowbreak\GL_n(\A)/\GL_n(\sO)$ 上一个确定的函数 $f$，它唯一到相差一个标量。因此，几何上我们要构造 $D^b(\sBun_n)$ 中一个特定的对象；当 $k=\Fq$ 时，它给出这个函数。

对一般的尖点自守表示 $\pi$，很难在 $\pi\subset\Funct_{\cusp}(\GL_n(\cK)\backslash\allowbreak\GL_n(\A))$ 中典范地选出一个函数。\footnote{由于我们的群是 $\GL_n$，这样的选择存在，即所谓的新形式，但从几何观点看并不那么自然。}因此，在相应的 $\sBun_n$ 版本上该构造什么样的层，并不清楚。

\section{附录：层-函数对应}

\obx{4.1}{ob:4.1}
这里回顾 Grothendieck 的层-函数对应的若干事实。

首先，若 $\mathcal{Y}$ 是（任意基域上的）概形，可以定义范畴 $D^b(\mathcal{Y})$，称为上同调可构造的 $\ell$-进层的导出范畴。

若 $\phi:\mathcal{Y}_1\to\mathcal{Y}_2$ 是概形间的态射，则有拉回与直接像函子
\[
\phi^*:D^b(\mathcal{Y}_2)\to D^b(\mathcal{Y}_1),\qquad
\phi_!:D^b(\mathcal{Y}_1)\to D^b(\mathcal{Y}_2).
\]

若 $\Fc_1$ 与 $\Fc_2$ 是 $D^b(\mathcal{Y})$ 的两个对象，可以定义它们的张量积 $\Fc_1\otimes\Fc_2\in D^b(\mathcal{Y})$。

\obx{4.2}{ob:4.2}
现设 $\mathcal{Y}$ 定义在基域 $\Fq$ 上。记 $Y$ 为 $\mathcal{Y}$ 的 $\Fq$-值点集（它是有限集），$\Funct(Y)$ 为其上 $\Ql$-值函数构成的向量空间。$D^b(\mathcal{Y})$ 的每个对象 $\Fc$ 都给出 $\Funct(Y)$ 中一个元素 $f$。

首先，若 $\mathcal{Y}$ 是点概形 $\Spec(\Fq)$，则 $D^b(\Spec(\Fq))$ 的对象就是有限维 $\Ql$-向量空间构成的有界复形，带 $\mathrm{Gal}(\Flq/\Fq)$ 的作用。对应的函数（此时只是 $\Ql$ 的一个元素）是 $\mathrm{Gal}(\Flq/\Fq)$ 中 Frobenius 元的迹的交替和。对一般的 $\mathcal{Y}$、$\Fc\in D^b(\mathcal{Y})$ 与 $y\in Y$，$f$ 在 $y$ 处的值等于上述规则赋给 $\Fc_y$ 的数，其中 $\Fc_y\in D^b(\Spec(\Fq))$ 是 $\Fc$ 沿对应于 $y$ 的态射 $\Spec(\Fq)\to\mathcal{Y}$ 的拉回。

对应 $\Fc\mapsto f$ 的基本性质如下。张量积 $\Fc_1\otimes\Fc_2$ 对应 $f_1\cdot f_2$（逐点乘积）。对映射 $\phi:\mathcal{Y}_1\to\mathcal{Y}_2$，拉回 $\phi^*(\Fc)$ 对应 $\phi^*(f)$，推前同理。（就证明而言，前两条性质由定义立得，第三条即 Grothendieck--Lefschetz 迹公式。）

\begin{annote}
第三条性质值得记住：它把「沿态射推前」与「逐纤维求和」等同起来。前文 \S2.5 中 $(h^{\rightarrow})_!$ 的求和、\S2.8 中 $q_!$ 的积分，都是它在函数侧的表现；到了层侧，同一件事由上同调的直接像承担，而迹公式保证两边算出的函数一致。\par\vspace{0.15em}{\raggedleft 【来源：原书 \S4.2（A）；Grothendieck--Lefschetz 迹公式的标准表述（B）】\par}
\end{annote}

\obx{4.3}{ob:4.3}
事实证明，许多情形下定义在 $\mathcal{Y}/\Fq$ 的 $Y$ 上的「有趣」函数，都按上述方式来自 $\ell$-进层。某些情形下也可以反向构造。下面是一个例子。

设 $A$ 是连通交换代数群。若对象 $\mathcal{X}\in D^b(A)$ 是一维局部系统，其在单位元处的纤维平凡，且 $m^*(\mathcal{X})\simeq\mathcal{X}\boxtimes\mathcal{X}$（其中 $m:A\times A\to A$ 是乘法映射），则称 $\mathcal{X}$ 为特征层。

此时，对应的有限群 $A=A(\Fq)$ 上的函数 $\chi$ 是特征。

反之，给定 $A$ 的特征 $\chi$，可如下重构 $\mathcal{X}$。考虑 Lang 同态 $L:A\to A$，
\[
a\mapsto\mathrm{Fr}(a)/a,
\]
其中 $\mathrm{Fr}:A\to A$ 是 Frobenius 自同态。则 $L$ 是 \'etale Galois 覆盖，其覆盖变换群等于 $\ker(L)=A$。

所求的特征层 $\mathcal{X}$，由 $\pi_1(A)$ 的如下一维表示所对应的 $\ell$-进层给出：
\[
\pi_1(A)\to\ker(L)=A\xrightarrow{\chi}\Ql^\times .
\]

把这一构造用于 $A$ 是仿射直线 $\mathbb{A}^1$ 的情形，得到 $\mathbb{A}^1$ 上称为 Artin--Schreier 层的 $\ell$-进层。类似地，$A$ 是乘法群 $\mathbb{G}_m$ 时得到其上的 Kummer 层。

另一个例子是 $A=\Pic^0(X)$：由此可恢复 Lang 对函数域上整体类域论的证明中所用的构造；参看 [Se]。

\par\bigskip
\section*{参考文献}
\addcontentsline{toc}{section}{参考文献}
\markright{参考文献}

{\small
\noindent\textit{背景读物}：除下列研究论文（见后）外，背景材料另有如下来源。\par
\smallskip
\noindent 代数几何（概形、凝聚层等）：Hartshorne, \emph{Algebraic Geometry}；Kempf, \emph{Algerbraic Varieties}；Mumford, \emph{The Red Book of Varieties and Schemes}。\par
\smallskip
\noindent 局部环与局部域、阿代尔、Frobenius 元素及其他代数数论概念：Cassels--Fr\"ohlich, \emph{Proceedings of the Instructional Conference on Algebraic Number Theory}；Serre, \emph{Corps locaux} 与 \emph{Groupes alg\'ebriques et corps de classes}。\par
\smallskip
\noindent $p$-进群的表示与自守函数理论导引：Gelfand--Graev--Piatetskii-Shapiro, \emph{Representation Theory and Automorphic Functions}；Gelbart, \emph{Automorphic Forms on Adele Groups}。关于 $p$-进群表示的高级课程（强烈推荐）见 [Be]。\par
\smallskip
\noindent \'Etale 层：Milne, \emph{Etale Cohomology}。\par
\smallskip
}

\begin{itemize}[leftmargin=4.6em,itemsep=0.22em,topsep=0.3em,label={}]
\item[{[Be]}] J. Bernstein, \emph{Lectures on Representations of $p$-Adic Groups}, preprint, School of Mathematical Sciences, Tel Aviv University, Tel Aviv, Israel; available online from \texttt{www.math.tau.ac.il/\~{}bernstei}.
\item[{[BD]}] A. Beilinson and V. Drinfeld, \emph{Quantization of Hitchin's Integrable System and Hecke Eigen-Sheaves}, preprint; available online from \texttt{www.math.uchicago.edu/\~{}benzvi}.
\item[{[De]}] P. Deligne, La conjecture de Weil II, \emph{Publ.\ IHES}, \textbf{52}(1981), 313--428.
\item[{[Dr]}] V. Drinfeld, Two-dimensional $\ell$-adic representations of the fundamental group of a curve over a finite field and automorphic forms on $GL(2)$, \emph{Amer.\ J.\ Math.}, \textbf{105}(1983), 85--114.
\item[{[L]}] G. Laumon, Correspondance de Langlands g\'eom\'etrique pour les corps de fonctions, \emph{Duke Math.\ J.}, \textbf{54}(1987), 309--359.
\item[{[L-MB]}] G. Laumon and L. Morret-Bailly, \emph{Champs alg\'ebriques}, Springer-Verlag, New York, 2000.
\item[{[Mi]}] J. Milne, \emph{Etale Cohomology}, Princeton Mathematical Series 33, Princeton University Press, Princeton, NJ, 1980.
\item[{[Se]}] J.-P. Serre, \emph{Algebraic Groups and Class Fields}, Graduate Texts in Mathematics 117, Springer-Verlag, New York, 1998.
\end{itemize}
